\newpage
\section*{Resumen}
\label{sec:resumen}
\addcontentsline{toc}{chapter}{Resumen}
El presente Trabajo Final aborda el diseño, la implementación y la validación metrológica de un sistema integral de software para la automatización de la captura y el análisis de ensayos de impulso tipo rayo (\qty[parse-numbers=false]{1.2/50}{\micro\second}), basado en los lineamientos de las normas IEC 60060-1~\cite{iec60060_1} e IEC 61083-2~\cite{iec61083_2}. El desarrollo responde a la necesidad de optimizar los procedimientos manuales de medición en el Laboratorio de Alta Tensión (FCEFyN-UNC), reduciendo los tiempos de ejecución y eliminando el error operacional humano.

El sistema integra una capa de abstracción de hardware basada en el protocolo VISA y comandos SCPI para controlar un osciloscopio digital GW Instek GDS-1102A-U~\cite{gds1000au_ProgrammingManual}. Desarrollado en Python bajo la arquitectura Modelo-Vista-Presentador (MVP), garantiza el desacoplamiento entre la interfaz gráfica y la lógica funcional con sondeo asíncrono no bloqueante.

El motor de cálculo numérico implementa el algoritmo establecido en la IEC 60060-1. Emplea regresión no lineal mediante el algoritmo de Levenberg-Marquardt sobre una función de doble exponencial, y aplica un filtro digital de fase cero para atenuar el ruido, lo que permite sintetizar la curva de tensión de ensayo~\cite{iec60060_1}. 

La gestión de la información se administra jerárquicamente en formato HDF5, integrando respaldos en texto plano (CSV) y exportación a planillas de cálculo.

La calibración del software se realizó procesando los registros patrón del generador de datos de prueba (TDG) de la norma IEC 61083-2~\cite{iec61083_2}, para verificar si los errores en el cálculo de los parámetros característicos ($U_t$, $T_1$, $T_2$, $T_c$ y $\beta'$) cumplen las tolerancias admisibles.

\section*{Área Temática}
Computación e Informática, Digitales.

\section*{Asignaturas}
Informática, Métodos numéricos, Probabilidad y estadística.

\section*{Palabras Claves}
Ensayo de impulso atmosférico, alta tensión, procesamiento digital de señales, calibración de software, normas IEC 60060-1 e IEC 61083-2

\newpage
\section*{Abstract}
This Final Project addresses the design, implementation, and metrological validation of a comprehensive software system for the automation of the acquisition and analysis of lightning impulse tests (\qty[parse-numbers=false]{1.2/50}{\micro\second}), based on the guidelines of the IEC 60060-1~\cite{iec60060_1} and IEC 61083-2~\cite{iec61083_2} standards. The development responds to the need to optimize manual measurement procedures at the High Voltage Laboratory (FCEFyN-UNC), reducing execution times and eliminating human operational error.

The system integrates a hardware abstraction layer based on the VISA protocol and SCPI commands to control a GW Instek GDS-1102A-U digital oscilloscope~\cite{gds1000au_ProgrammingManual}. Developed in Python under the Model-View-Presenter (MVP) architecture, it ensures the decoupling between the graphical user interface and the functional logic with non-blocking asynchronous polling.

The numerical calculation engine implements the algorithm established in IEC 60060-1. It employs non-linear regression via the Levenberg-Marquardt algorithm on a double exponential function, and applies a zero-phase digital filter to attenuate noise, enabling the synthesis of the test voltage curve~\cite{iec60060_1}.

Information management is hierarchically administered in HDF5 format, integrating plain text backups (CSV) and spreadsheet exports.

Software calibration was performed by processing the reference records from the Test Data Generator (TDG) of the IEC 61083-2 standard~\cite{iec61083_2}, to verify whether the errors in the calculation of the characteristic parameters ($U_t$, $T_1$, $T_2$, $T_c$, and $\beta'$) comply with the allowable tolerances.
\section*{Thematic area}
Computer Science and Informatics, Digital Systems.

\section*{Subjects}
Informatics, Numerical Methods, Probability and Statistics.

\section*{Key Words}
Lightning impulse test, high voltage, digital signal processing, software calibration, IEC 60060-1 and IEC 61083-2 standards.

\newpage
\section*{Resumo}
Este Trabalho Final aborda o projeto, a implementação e a validação metrológica de um sistema integral de software para a automação da captura e análise de ensaios de impulso atmosférico (\qty[parse-numbers=false]{1.2/50}{\micro\second}), baseado nas diretrizes das normas IEC 60060-1~\cite{iec60060_1} e IEC 61083-2~\cite{iec61083_2}. O desenvolvimento atende à necessidade de otimizar os procedimentos manuais de medição no Laboratório de Alta Tensão (FCEFyN-UNC), reduzindo os tempos de execução e eliminando o erro operacional humano.

O sistema integra uma camada de abstração de hardware baseada no protocolo VISA e comandos SCPI para controlar um osciloscópio digital GW Instek GDS-1102A-U~\cite{gds1000au_ProgrammingManual}. Desenvolvido em Python sob a arquitetura Modelo-Visão-Apresentador (MVP), garante o desacoplamento entre a interface gráfica e a lógica funcional com sondagem assíncrona não bloqueante.

O motor de cálculo numérico implementa o algoritmo estabelecido na IEC 60060-1. Emprega regressão não linear por meio do algoritmo de Levenberg-Marquardt sobre uma função dupla exponencial e aplica um filtro digital de fase zero para atenuar o ruído, o que permite sintetizar a curva de tensão de ensaio~\cite{iec60060_1}.

O gerenciamento de dados é administrado hierarquicamente no formato HDF5, integrando backups em texto simples (CSV) e exportação para planilhas de cálculo.

A calibração do software foi realizada processando os registros de referência do gerador de dados de teste (TDG) da norma IEC 61083-2~\cite{iec61083_2}, para verificar se os erros no cálculo dos parâmetros característicos ($U_t$, $T_1$, $T_2$, $T_c$ e $\beta'$) cumprem as tolerâncias admissíveis.

\section*{Área Temática}
Ciência da Computação e Informática, Sistemas Digitais.

\section*{Assuntos}
Informática, Métodos Numéricos, Probabilidade e Estatística.

\section*{Palavras Chaves}
Ensaio de impulso atmosférico, alta tensão, processamento digital de sinais, calibração de software, normas IEC 60060-1 e IEC 61083-2.